\documentclass[11pt,a4paper]{article}

\usepackage[a4paper,margin=1in]{geometry}
\usepackage{fontspec}
\IfFontExistsTF{Times New Roman}
  {\setmainfont{Times New Roman}}
  {\IfFontExistsTF{TeX Gyre Termes}{\setmainfont{TeX Gyre Termes}}{\setmainfont{Times}}}
\usepackage{newtxmath}
\usepackage{amsmath}
\usepackage{booktabs}
\usepackage{longtable}
\usepackage{array}
\usepackage{enumitem}
\usepackage{xcolor}
\usepackage{listings}
\usepackage{graphicx}
\usepackage{float}
\usepackage[hidelinks]{hyperref}
\usepackage{xurl}
\usepackage{fancyhdr}
\usepackage{microtype}

\setlength{\parindent}{0pt}
\setlength{\parskip}{0.65em}
\setlength{\emergencystretch}{3em}
\setlist[itemize]{leftmargin=1.4em,itemsep=0.2em,topsep=0.2em}
\setlist[enumerate]{leftmargin=1.6em,itemsep=0.2em,topsep=0.2em}

\pagestyle{fancy}
\fancyhf{}
\lhead{\small Ticket 10 Study Report}
\rhead{\small Delta and Gamma Diagnostics}
\cfoot{\thepage}
\renewcommand{\headrulewidth}{0.3pt}

\lstset{
  basicstyle=\ttfamily\small,
  breaklines=true,
  frame=single,
  columns=fullflexible,
  keepspaces=true,
  rulecolor=\color{black!30},
  xleftmargin=0.5em,
  xrightmargin=0.5em
}

\newcommand{\code}[1]{\texttt{#1}}
\newcommand{\file}[1]{\path{#1}}
\newcommand{\func}[1]{\path{#1}}
\newcolumntype{L}[1]{>{\raggedright\arraybackslash}p{#1}}
\newcolumntype{P}[1]{>{\raggedright\arraybackslash\ttfamily\footnotesize}p{#1}}

\title{\textbf{Ticket 10 Study Report: Delta and Gamma Diagnostics}\\
\large SURF2026 American Option Risk Surfaces}
\author{Codex-assisted implementation report}
\date{June 16, 2026}

\begin{document}
\maketitle

\begin{abstract}
Ticket 10 adds numerical Delta and Gamma diagnostics for American CN/PSOR value grids. It is
not a new pricing algorithm. It reads validated American option value grids from the Ticket 04
solver, uses the Ticket 09 boundary-extraction layer for diagnostic context, computes central
finite-difference Greeks, and separates ordinary interior nodes from fragile regions near
payoff kinks, maturity, and extracted free-boundary bands. This report is written for beginner
researchers with accounting, Python, data, and machine-learning experience who are still learning
option Greeks and finite-difference instability.
\end{abstract}

\tableofcontents
\newpage

\section{Role of Ticket 10 in the Whole Solver Project}

Tickets 01 through 09 built the validation ladder before Greeks:
\begin{itemize}
  \item Ticket 01 created payoff and Black-Scholes closed-form utilities.
  \item Tickets 02 and 03 built and validated the finite-difference European machinery.
  \item Ticket 04 implemented the American CN/PSOR obstacle core.
  \item Tickets 05 through 07 validated American put, no-dividend call, and dividend-paying call behavior.
  \item Ticket 08 diagnosed obstacle and complementarity conditions.
  \item Ticket 09 extracted approximate free-boundary curves from the continuation premium \(U-\Phi\).
\end{itemize}

Ticket 10 is deliberately after that sequence. Delta and Gamma are derivatives of the computed
value surface. Derivatives amplify numerical noise, so the project first needed price validation,
LCP diagnostics, and boundary context. Ticket 10 therefore asks: once the value grid is credible,
what can we safely say about finite-difference Delta and Gamma, and where should we be cautious?

\section{Theory and Methodology Connection}

The documented source basis for this ticket is internal:
\begin{itemize}
  \item the literature map's warning that Greeks are delicate near payoff kinks and free boundaries;
  \item the formulation note's definitions of Delta and Gamma;
  \item the solver validation plan's requirement to report Delta bounds, Gamma fragility, kink masks, and boundary masks;
  \item prior Ticket 04--09 reports, which establish the American solver, diagnostics, and boundary extraction.
\end{itemize}

No BibTeX file is created for Ticket 10. The report does not invent external bibliography.
It uses the project documents as the source basis and treats Greek diagnostics as part of the
larger American-option validation ladder.

\section{Delta and Gamma in Plain Language}

Delta measures how much the option value changes when the stock price changes a little:
\[
\Delta(S,\tau) = \frac{\partial U}{\partial S}(S,\tau).
\]

Gamma measures the curvature of the option value with respect to the stock price:
\[
\Gamma(S,\tau) = \frac{\partial^2 U}{\partial S^2}(S,\tau).
\]

For intuition:
\begin{itemize}
  \item a put Delta is usually between \(-1\) and \(0\);
  \item a call Delta is usually between \(0\) and \(1\);
  \item Gamma is usually more fragile than price or Delta because it is a second derivative.
\end{itemize}

In this ticket, Delta and Gamma are numerical diagnostics, not exact analytical Greeks. The
project does not yet use them as machine-learning labels or production risk measures.

\section{Finite-Difference Greek Formulas}

On a uniform spot grid \(S_i=i\Delta S\), Ticket 10 uses central differences at interior nodes.
For a value grid \(U_i^n \approx U(S_i,\tau_n)\), the Delta approximation is
\[
\Delta_i^n \approx \frac{U_{i+1}^n - U_{i-1}^n}{2\Delta S}.
\]

The Gamma approximation is
\[
\Gamma_i^n \approx \frac{U_{i+1}^n - 2U_i^n + U_{i-1}^n}{(\Delta S)^2}.
\]

Boundary nodes are reported as \(\mathrm{NaN}\). This is intentional. One-sided endpoint
derivatives are possible, but they are not as comparable to the central interior formulas. Ticket
10 chooses clarity over pretending endpoint Greeks are equally reliable.

\begin{longtable}{P{0.25\textwidth}L{0.28\textwidth}L{0.35\textwidth}}
\toprule
Quantity & Formula & Meaning \\
\midrule
\endhead
\path{Delta} & \(\left(U_{i+1}-U_{i-1}\right)/(2\Delta S)\) & First derivative; approximate spot sensitivity. \\
\path{Gamma} & \(\left(U_{i+1}-2U_i+U_{i-1}\right)/(\Delta S)^2\) & Second derivative; approximate curvature. \\
\path{Boundary nodes} & \(\mathrm{NaN}\) & Endpoint Greeks are excluded from headline diagnostics. \\
\path{Strict mask} & finite Greeks away from caution bands & Used for cleaner interior summaries. \\
\bottomrule
\end{longtable}

\section{Why Greeks Come After Boundary Extraction}

The American exercise boundary is a non-smooth feature of the value surface. Near that transition,
small changes in spot can move a node from continuation-like behavior to exercise-like behavior.
That is exactly where Gamma can spike or become visually noisy.

Ticket 09 extracted approximate boundary curves. Ticket 10 uses those curves only as diagnostic
context. It does not claim the boundary is exact. It marks nearby nodes so the report can compare:
\begin{itemize}
  \item full-grid Greek summaries;
  \item away-from-boundary summaries;
  \item strict-interior summaries that also remove payoff-kink and maturity rows.
\end{itemize}

\section{Boundary-Near Instability}

Boundary-near instability means the Greek estimate is influenced by the discrete exercise
transition. In an American put, the transition typically appears on the low-spot side. In a
dividend-paying American call, it may appear on the high-spot side. Ticket 10 uses the boundary
points from Ticket 09 and marks nodes within two grid steps of each found boundary.

This mask is not boundary extraction. Boundary extraction happened in Ticket 09. Ticket 10 uses
the extracted boundary as a caution label for Greeks.

\section{Payoff Kink and Maturity-Row Caution}

At maturity, the option value equals payoff. Payoffs have kinks at the strike \(K\). A kink is a
point where the slope changes abruptly. Delta can jump there, and Gamma can look like a large
spike on a discrete grid.

Ticket 10 marks:
\begin{itemize}
  \item nodes within two grid steps of \(S=K\);
  \item the maturity row \(\tau=0\);
  \item nodes within two grid steps of extracted boundary points.
\end{itemize}

\begin{longtable}{P{0.30\textwidth}L{0.50\textwidth}}
\toprule
Caution region & Why it matters \\
\midrule
\endhead
\path{maturity row} & The value equals payoff, so the kink is strongest and Gamma can be large. \\
\path{payoff-kink band} & Central differences straddle a non-smooth payoff slope change near \(S=K\). \\
\path{boundary-near band} & Central differences may straddle exercise-like and continuation-like behavior. \\
\path{domain boundary} & Endpoint formulas are excluded and reported as \(\mathrm{NaN}\). \\
\bottomrule
\end{longtable}

\section{Implementation Design Decisions}

\begin{longtable}{P{0.35\textwidth}L{0.47\textwidth}}
\toprule
Decision & Reason \\
\midrule
\endhead
\path{src/american_risk_surfaces/diagnostics/greeks.py} & Keeps Greek diagnostics separate from the solver core. \\
\path{central differences only} & Matches the uniform-grid finite-difference setup used in earlier tickets. \\
\path{NaN boundary Greeks} & Avoids overstating endpoint derivative reliability. \\
\path{Ticket 09 boundary mask} & Reuses existing boundary extraction without changing pricing algorithms. \\
\path{strict interior mask} & Separates headline summaries from kink, boundary, maturity, and nonfinite nodes. \\
\path{no Rannacher smoothing} & Smoothing is explicitly deferred to a later ticket. \\
\bottomrule
\end{longtable}

\section{Code-Level Function Explanations}

\begin{longtable}{P{0.35\textwidth}L{0.47\textwidth}}
\toprule
Function or class & Role \\
\midrule
\endhead
\func{GreekArrays} & Stores Delta, Gamma, and finite masks. \\
\func{GreekDiagnosticMasks} & Stores boundary-near, kink-near, maturity-row, and strict-interior masks. \\
\func{GreekTimeDiagnosticRow} & Records by-time counts and Greek warnings. \\
\func{GreekDiagnosticSummary} & Records whole-case summary metrics. \\
\func{GreekDiagnostics} & Bundles arrays, masks, summary, and by-time rows. \\
\func{finite_difference_delta} & Computes central Delta and sets boundary nodes to \(\mathrm{NaN}\). \\
\func{finite_difference_gamma} & Computes central Gamma and sets boundary nodes to \(\mathrm{NaN}\). \\
\func{payoff_kink_mask} & Marks nodes near \(S=K\). \\
\func{boundary_near_mask} & Marks nodes near Ticket 09 extracted boundary points. \\
\func{maturity_row_mask} & Marks the row where \(\tau=0\). \\
\func{diagnose_greek_result} & Main API for producing Greek diagnostics from an American CN/PSOR result. \\
\func{selected_greek_profile_rows} & Produces selected moneyness and tau rows for CSV/report review. \\
\func{greek_by_time_rows} & Exposes the by-time diagnostic rows. \\
\bottomrule
\end{longtable}

\section{Diagnostic Cases and Parameters}

{\small
\setlength{\tabcolsep}{3pt}
\begin{longtable}{P{0.22\textwidth}P{0.08\textwidth}P{0.06\textwidth}P{0.06\textwidth}P{0.06\textwidth}P{0.06\textwidth}P{0.08\textwidth}P{0.08\textwidth}}
\toprule
Case & Type & \(K\) & \(T\) & \(r\) & \(q\) & \(\sigma\) & \(M=N\) \\
\midrule
\endhead
\path{american_put_medium} & put & 1.0 & 1.0 & 0.05 & 0.02 & 0.20 & 120 \\
\path{dividend_call_medium} & call & 1.0 & 1.0 & 0.05 & 0.08 & 0.20 & 120 \\
\bottomrule
\end{longtable}
}

Both cases use \(S_{\max}=4K\). The Greek report is not a grid-sensitivity study. Later tickets
can vary grid size, domain size, smoothing settings, and dataset construction.

\section{Greek Result Tables and Interpretation}

The summary CSV contains two rows. The headline values below are from
\file{results/01_solver_validation/tables/ticket_10_greek_diagnostics_summary.csv}.

{\small
\setlength{\tabcolsep}{3pt}
\begin{longtable}{P{0.21\textwidth}P{0.07\textwidth}P{0.10\textwidth}P{0.10\textwidth}P{0.10\textwidth}P{0.12\textwidth}P{0.16\textwidth}}
\toprule
Case & Type & finite \(\Delta\) & finite \(\Gamma\) & max \(|\Gamma|\) & strict max \(|\Gamma|\) & Status \\
\midrule
\endhead
\path{american_put_medium} & put & 14399 & 14399 & 30.0000 & 3.1800 & \path{PASS_WITH_CAUTIONS} \\
\path{dividend_call_medium} & call & 14399 & 14399 & 30.0000 & 2.6799 & \path{PASS_WITH_CAUTIONS} \\
\bottomrule
\end{longtable}
}

The full-grid maximum \(|\Gamma|\) is \(30\) for both cases. This is not surprising: the
maturity payoff kink creates a large second-difference response. After excluding maturity, kink,
boundary-near, and endpoint nodes, the strict maximum \(|\Gamma|\) is much smaller. This is the
main Ticket 10 lesson: do not hide the spike, but do not let it dominate the interpretation of
stable interior Greeks either.

\begin{longtable}{P{0.23\textwidth}P{0.14\textwidth}P{0.14\textwidth}P{0.14\textwidth}P{0.14\textwidth}}
\toprule
Case & boundary-near nodes & kink-near nodes & strict nodes & strict violations \\
\midrule
\endhead
\path{american_put_medium} & 480 & 605 & 13243 & 0 lower, 0 upper, 0 negative \(\Gamma\) \\
\path{dividend_call_medium} & 480 & 605 & 13231 & 0 lower, 0 upper, 0 negative \(\Gamma\) \\
\bottomrule
\end{longtable}

Selected profile rows show the caution masks in action. For example, near the strike at early
time, the put has \(\Gamma \approx 22.63\) at \(S/K=1.0\), and that row is marked as both
boundary-near and kink-near. At \(S/K=1.2\), the same early-time put row has a tiny value and
small Gamma and is marked strict-interior. For the dividend call, the strike row is kink-near,
while the high-spot exercise-side region can be boundary-near.

\section{Figure Interpretation}

\begin{figure}[H]
  \centering
  \includegraphics[width=0.92\textwidth]{results/01_solver_validation/figures/ticket_10_american_put_delta_gamma_profiles.png}
  \caption{American put Delta and Gamma profiles at selected time-to-maturity levels.}
\end{figure}

The put figure shows Delta mostly between \(-1\) and \(0\). Gamma is visibly largest near the
payoff kink and transition regions, which is why the report uses caution masks.

\begin{figure}[H]
  \centering
  \includegraphics[width=0.92\textwidth]{results/01_solver_validation/figures/ticket_10_dividend_call_delta_gamma_profiles.png}
  \caption{Dividend-paying American call Delta and Gamma profiles at selected time-to-maturity levels.}
\end{figure}

The call figure shows Delta mostly between \(0\) and \(1\). Gamma again highlights curvature
near the strike and transition behavior.

\begin{figure}[H]
  \centering
  \includegraphics[width=0.92\textwidth]{results/01_solver_validation/figures/ticket_10_gamma_boundary_diagnostic.png}
  \caption{Near-boundary versus away-from-boundary maximum Gamma magnitude.}
\end{figure}

This figure is a diagnostic comparison, not a proof of exact boundary accuracy. The boundary band
comes from Ticket 09 and is threshold-based.

\begin{figure}[H]
  \centering
  \includegraphics[width=0.92\textwidth]{results/01_solver_validation/figures/ticket_10_gamma_by_time.png}
  \caption{Full-grid versus strict-mask maximum absolute Gamma over time.}
\end{figure}

The full-grid curve includes maturity and kink effects. The strict curve removes the main caution
regions and better represents stable interior behavior.

\section{Testing Strategy}

\begin{longtable}{P{0.34\textwidth}L{0.45\textwidth}}
\toprule
Test category & What it proves \\
\midrule
\endhead
Synthetic constant, linear, quadratic functions & Central differences reproduce known derivatives on a uniform grid. \\
Shape and boundary tests & Delta and Gamma match value-grid shape and use \(\mathrm{NaN}\) at boundaries. \\
Invalid input tests & Bad grids, bad value shapes, and bad diagnostic arguments raise \path{ValueError}. \\
Mask tests & Kink, maturity, and boundary-near masks behave predictably. \\
Real American put and dividend call tests & Diagnostics produce finite arrays and summaries on solver outputs. \\
CSV and figure tests & The experiment writes stable artifacts with expected columns and nonempty figure files. \\
Scope tests & No Rannacher smoothing, stress maps, datasets, or neural-network APIs are introduced. \\
\bottomrule
\end{longtable}

\section{Test Results and Interpretation}

The final test command is:
\begin{lstlisting}
PYTHONPYCACHEPREFIX=/private/tmp/codex_pycache PYTHONPATH=src MPLCONFIGDIR=/private/tmp/matplotlib-cache FC_CACHEDIR=/private/tmp/fontconfig-cache python3 -m unittest discover -s tests
\end{lstlisting}

The observed result after implementation was:
\begin{lstlisting}
Ran 121 tests in 2.198s
OK
\end{lstlisting}

This means the new Ticket 10 tests passed along with prior tickets' tests in the current
environment. It does not mean Gamma is analytically exact, and it does not validate later
Rannacher smoothing, stress-map generation, dataset labels, or neural surrogate modeling.

\section{Implementation Pitfalls and Debugging Notes}

{\small
\setlength{\tabcolsep}{3pt}
\begin{longtable}{P{0.25\textwidth}L{0.23\textwidth}L{0.23\textwidth}L{0.17\textwidth}}
\toprule
Event or trap & Symptom & Fix or mitigation & Beginner lesson \\
\midrule
\endhead
Initial TDD red run & 11 errors from missing \path{greeks.py} and \path{07_greek_diagnostics.py}. & Implemented the diagnostic module and experiment, then reran the suite. & A red test should fail for the expected missing behavior. \\
Endpoint derivative misuse & Boundary Greeks can look precise when they are one-sided artifacts. & Boundary Delta/Gamma are reported as \(\mathrm{NaN}\). & Not every array cell deserves a Greek. \\
Overinterpreting Gamma spikes & Full-grid max \(|\Gamma|\) reaches 30 near kink/maturity. & Report full and strict-mask Gamma separately. & Gamma spikes are information, not automatically bugs. \\
Treating boundary estimates as exact & Boundary masks inherit Ticket 09 threshold choices. & Report them as caution context only. & A diagnostic mask is not analytical truth. \\
Adding smoothing too early & Rannacher smoothing could reduce noise but change the experiment. & Deferred smoothing to a later named ticket. & Do not silently change the method to make plots nicer. \\
Using Greeks as dataset labels too early & Diagnostic Greeks are tempting labels. & Report that dataset work is out of scope. & Validation comes before labels. \\
\bottomrule
\end{longtable}
}

\section{Limitations}

\begin{longtable}{P{0.32\textwidth}L{0.50\textwidth}}
\toprule
Limitation & Meaning \\
\midrule
\endhead
Representative cases only & The experiment uses one American put and one dividend-paying American call setup. \\
No grid sensitivity & \(M=N=120\) is not a refinement study. \\
No Rannacher smoothing & Gamma noise is reported, not smoothed away. \\
Boundary masks are threshold-based & They inherit Ticket 09's approximate extraction logic. \\
No dataset labels & Greeks are not exported as machine-learning training labels. \\
No stress maps or neural work & Those belong to later project stages. \\
\bottomrule
\end{longtable}

\section{Lessons Learned / Study Notes}

\begin{itemize}
  \item Delta is usually easier to interpret than Gamma because it is a first derivative.
  \item Gamma is sensitive to non-smooth value features: payoff kink, maturity row, and exercise boundary.
  \item A useful Greek diagnostic should show both full-grid behavior and masked interior behavior.
  \item Masking is not hiding. It is labeling where the numerical derivative is fragile.
  \item A large Gamma near \(S=K\) at maturity is expected for a kinked payoff.
  \item The right response to noisy Gamma is careful reporting first, not immediate smoothing.
\end{itemize}

\section{Link to the Next Ticket}

The next ticket should not treat these Greeks as final labels yet. A natural next step is the
Rannacher smoothing comparison gate, where smoothed and unsmoothed prices, boundaries, Delta, and
Gamma can be compared explicitly. After that, grid/domain sensitivity and dataset construction
can be considered with clearer evidence.

\end{document}
